\section{A depth budget for the actual full source matrix}
\label{sec:weighted-cluster-depth-sufficiency}

\paragraph{Status.}
The complete author and two independent whole-proof readings
have passed; their exact check scopes are recorded separately.
No arithmetic validity of the stated cluster geometry is asserted.

The whole-product method of MA-MO-030 permits a more precise
estimate than a fixed bound on the total local multiplicity.
The polynomial weight pays for part of the pole order; positive
inward depth controls the rest.  Actual boundary factors require
a separate integrability condition, retained explicitly below.

Use the two source settings of (MA-MO.700): either a nontrivial
finite-order Hecke character with its dual and structural $S$,
or the rational principal sector with its actual polar completion.
Keep the original source Fourier maps, real smooth step $\phi$,
primitive finite functions $A_\omega$, and inherited norms.
Fix $0<a<1/2$, $s\ge0$, and put
\[
 r=s/2,\qquad p=r+1/2=(s+1)/2,\qquad
 0<\delta<\tfrac18\min\{a,1/2-a\},\qquad0<\ell\le1.
 \tag{MA-MO.1400}
\]
Fix also $0<\epsilon_0<\delta/8$ for the right-boundary limits.
Assume that every completed centered zero lies in $|\Re z|\le a$,
for both dual sectors when applicable.  At $t_n=n\ell$, let
$Z_{\omega,n}$ be the full multiset in
\[
 a-\delta\le\Re\lambda\le a,
                  \qquad |\Im\lambda-t_n|\le2\ell.
 \tag{MA-MO.1401}
\]
Write $k_{\omega,n}$ for its total multiplicity and
$m_{\omega,n}$ for the total multiplicity of its factors on
$\Re\lambda=a$.  Impose
\[
                         m_{\omega,n}<p
                          \qquad(\omega,n).
 \tag{MA-MO.1402}
\]
This concerns the total boundary multiplicity in the rectangle,
even when the boundary zeros have distinct ordinates.
It is stronger than merely retaining each of those factors
individually.  When positive-depth factors exist, define
\[
 d_{\omega,n}
   =\min_{\substack{\lambda\in Z_{\omega,n}\\\Re\lambda<a}}
                               (a-\Re\lambda)\in(0,\delta].
 \tag{MA-MO.1403}
\]
If no such factor exists, put $d_{\omega,n}=1$.
In that case $k_{\omega,n}=m_{\omega,n}<p$.
Define the explicit cost
\[
 B_{\omega,n}=
 \begin{cases}
 1,&k_{\omega,n}<p,\\
 \bigl(1+\log(1/d_{\omega,n})\bigr)^{1/2},&k_{\omega,n}=p,\\
 d_{\omega,n}^{-(k_{\omega,n}-p)},&k_{\omega,n}>p.
 \end{cases}
 \tag{MA-MO.1404}
\]
In the last two cases (MA-MO.1402) guarantees the presence of
a positive-depth factor.  Suppose that, for fixed $C,D$,
\[
                  B_{\omega,n}\le C(2+|n|)^D
                              \qquad(\omega,n).
 \tag{MA-MO.1405}
\]
All multiplicities in these definitions are the actual ones.

\paragraph{MA-MO-035 (a polynomial depth budget suffices).}
Under (MA-MO.1401)--(MA-MO.1405) and the closed-strip premise,
the original $C_g(P_+-P_-)$ extends uniquely in trace class
on $H_{\Sigma,a,s}$ in the nontrivial sector, and on
$\mathcal W_{a,s}=H_{\Sigma,a,s}\oplus\mathbb C^2$ in the
rational principal sector, for every compact smooth $g$.
The extensions agree with the original wide-inverse source maps
and depend continuously on $g$ in trace norm on each fixed
compact support.

All completed jets are retained.  The actual relative-trace,
Fredholm-determinant and prime formulas of MA-MO-029 apply,
with their original signs, finite polar correction and condition
$g(0)=0$ for the stated prime formulas.  The assertion is a
sufficient condition; neither (MA-MO.1402) nor (MA-MO.1405)
is claimed necessary.

\begin{proof}
\emph{1. A uniform pole-product estimate.}
Consider one rectangle and suppress its indices.  Enumerate
its $k$ factors by
$\lambda_j=a-d_j+i(t_n+\gamma_j)$, with
$0\le d_j\le\delta$ and $|\gamma_j|\le2\ell$.
For $\epsilon>0$ let
\[
 R_\epsilon(v)=\prod_{j=1}^k
                  (\epsilon+d_j+i(v-\gamma_j))^{-1}.
 \tag{MA-MO.1406}
\]
If $k\ge1$, the actual Laplace formula is
\[
 R_\epsilon(v)=\int_0^\infty e^{-ivt}e^{-\epsilon t}h(t)\,dt,
 \qquad
 h=\mathop{*}_{j=1}^k
       \bigl(\mathbf1_{t>0}e^{-d_jt}e^{i\gamma_jt}\bigr).
 \tag{MA-MO.1407}
\]
Consequently
$|h(t)|\le t^{k-1}/(k-1)!$.  All estimates below are
independent of the ordinates and their separations.
With the unitary plus Fourier transform,
\[
 \|R_\epsilon\|_{H^{-r}}^2
       =2\pi\int_0^\infty
           (1+t^2)^{-r}e^{-2\epsilon t}|h(t)|^2\,dt.
 \tag{MA-MO.1408}
\]

If $1\le k<p$, the simplex bound is integrable in (MA-MO.1408),
uniformly even when every depth is zero.  There are only finitely
many such integers $k$, so their bounds are controlled by a
constant depending on $s$ alone.

Suppose first that every depth is positive, with least depth $d$.
Then $|h(t)|\le e^{-dt}t^{k-1}/(k-1)!$.
If $k>p$, split the integral at one and use
$(1+t^2)^{-r}\le t^{-2r}$ for $t\ge1$.  With
$q_k=2(k-p)>0$, the result is bounded by
\[
 C+\frac{2\pi\Gamma(q_k)}{((k-1)!)^2}
                                      (2d)^{-q_k}.
 \tag{MA-MO.1409}
\]
Here the constant from $t\le1$ is uniform in $k$.
The gamma ratio has a uniform estimate, not a separate
constant for each growing multiplicity.  If $q_k\ge2$,
log-convexity of the gamma integral, between $2$ and
$2k-1$, gives
\[
 \Gamma(q_k)\le\Gamma(2k-1)=(2k-2)!,\qquad
 \frac{(2k-2)!}{((k-1)!)^2}\le4^{k-1}.
 \tag{MA-MO.1410}
\]
For $0<q_k<2$ there are only finitely many eligible integers
$k$ for this fixed $s$, and a constant depending on $s$ covers
them.  Since $q_k=2(k-p)$, the power of two in (MA-MO.1409)
cancels $4^{k-1}$ up to the fixed factor $4^{p-1}$.
It follows that $\|R_\epsilon\|_{H^{-r}}\le C_s d^{-(k-p)}$.
If $k=p$, which can occur only for this fixed integer $p$,
split instead at $1$ and $1/d$.  The middle integral is
bounded by $\log(1/d)$, and the tail by
$\int_1^\infty e^{-2u}\,du/u$.  This gives the critical
square-root logarithm in (MA-MO.1404).

Now allow $m\ge1$ boundary factors, where $m<p$, and put
$l=k-m$ for the number of positive-depth factors.
If $l=0$, then $k<p$ and the first case already applies.
Otherwise let $d$ be the least of their positive depths.
Convolving the separate simplex bounds gives both estimates
\[
 |h(t)|\le
 \min\left\{\frac{t^{k-1}}{(k-1)!},
             \frac{t^{m-1}d^{-l}}{(m-1)!}\right\}.
 \tag{MA-MO.1411}
\]
For the second, write the convolution parameter of the
positive-depth factors as $u$, bound $(t-u)^{m-1}$ by
$t^{m-1}$, and integrate
$e^{-du}u^{l-1}/(l-1)!$ over $u>0$ to get $d^{-l}$.
No exponential damping is assigned to the boundary factors.

Use the first bound up to $t=1/d$ and the second thereafter.
For $k>p$ the two large-$t$ integrals are bounded, apart from
their factorials, by
\[
 \int_1^{1/d}t^{2(k-p)-1}\,dt
       \le\frac{d^{-2(k-p)}}{2(k-p)},\qquad
 d^{-2l}\int_{1/d}^\infty t^{2(m-p)-1}\,dt
       =\frac{d^{-2(k-p)}}{2(p-m)}.
 \tag{MA-MO.1412}
\]
The first positive denominator is bounded away from zero
as $k$ varies through integers greater than the fixed $p$.
The second is also bounded away from zero, because $m$ ranges
over the finite set of positive integers less than $p$.
The factorial denominators only improve the bound.
For $k=p$, the first integral is $\log(1/d)$ and the second
is a fixed constant.  The case $k<p$ was already handled.
Thus in every nonempty case,
\[
 \sup_{0<\epsilon\le\epsilon_0}\|R_\epsilon\|_{H^{-r}}
          \le C_s B,\qquad
 R_\epsilon\longrightarrow R_0\quad\hbox{in }H^{-r}
                              \quad(\epsilon\downarrow0).
 \tag{MA-MO.1413}
\]
Convergence follows by dominated convergence in (MA-MO.1408)
for this fixed multiset.  A common modulus across all rectangles
is neither asserted nor needed.
For $k=0$ the unlocalized constant $R_\epsilon=1$ is not
put in $H^{-r}$; its compactly localized smooth term below
is handled separately.

\emph{2. The canceled holomorphic factors remain polynomial.}
Define $P_{\omega,n}$, $b_{\omega,n}=P_{\omega,n}/A_\omega$
and the nested rectangles exactly as in (MA-MO.903).
The ordinary local zero bound gives
\[
                    k_{\omega,n}=O(\log(2+|n|)).
 \tag{MA-MO.1414}
\]
The proof of (MA-MO.904) still gives a polynomial bound
for every fixed derivative of $b_{\omega,n}$, uniformly on
the smaller rectangle.  To check the dependence on $k$, remove
all finite-function zeros in the fixed disk centered at $2+it_n$.
The zero-free factor has polynomial inverse by the Euler lower
bound and Harnack.  For every zero in $Z_{\omega,n}$, the
numerator of its Blaschke factor is canceled by $P_{\omega,n}$;
the remaining factor costs a fixed constant.  All other disk
zeros are a fixed distance from the smaller rectangle, so each
also costs a fixed constant.  There are $O(\log(2+|n|))$
factors altogether.  Thus even a factor $C^{k_{\omega,n}}$
costs only a polynomial.  Fixed-radius Cauchy estimates give
the same conclusion for each fixed derivative.  The finite
Euler factors are invertible on this positive strip, and the
rational pole is outside the rectangle.

Use the fixed smooth partition $\chi_n$ of (MA-MO.910), and set
\[
 d_{\omega,n,\epsilon}(u)
   =\chi_n(u)b_{\omega,n}(a+\epsilon+iu)
                              R_{n,\epsilon}(u-t_n).
 \tag{MA-MO.1415}
\]
For $\epsilon>0$ this is precisely the localized actual
reciprocal.  Sobolev multiplication, (MA-MO.1405) and
(MA-MO.1413) yield, for some fixed $C,D_1$,
\[
 \|d_{\omega,n,\epsilon}\|_{H^{-r}}\le C(2+|n|)^{D_1},\qquad
 d_{\omega,n,\epsilon}\to d_{\omega,n,0}\quad\hbox{in }H^{-r}
                         \quad\text{for each fixed }n.
 \tag{MA-MO.1416}
\]
The estimates are uniform for $0\le\epsilon\le\epsilon_0$.
When $k=0$ the compactly supported smooth factor has this
bound directly.  This is the only use of the empty product.

The locally finite sum $j_{\omega,+}=\sum_n d_{\omega,n,0}$
is a tempered distribution.  Polynomial bounds and the rapid
decrease of localized Schwartz tests give
\[
 \frac1{A_\omega(a+\epsilon+i\,\cdot)}\to j_{\omega,+}
          \quad\hbox{in distributions},\qquad
                 A_\omega(a+i\,\cdot)j_{\omega,+}=1.
 \tag{MA-MO.1417}
\]
The second identity is taken after multiplying at $\epsilon>0$;
the forward multiplier and its derivatives have uniform polynomial
bounds.  For each fixed Schwartz $f$, the products
$j_{\omega,\epsilon}f$ also converge in $H^{-r}$.
These conclusions use dominated summation of (MA-MO.1416),
not a uniform-in-$n$ boundary-convergence modulus.

\emph{3. Actual source identification and trace-norm summation.}
For a two-variable Schwartz kernel $L$, the compact-distribution
Hermite lemma (MA-MO.715), translated to the $n$th window, gives
\[
 \|T_{d_{\omega,n,\epsilon},L}\|_{\mathcal S_1(H^r,H^R)}
     \le C_{L,N}(2+|n|)^{D_1-N}\qquad(R\ge0)
 \tag{MA-MO.1418}
\]
for every $N$, with finitely many Schwartz seminorms for fixed
$r,R,N$.  Choose $N>D_1+2$.  The series is uniformly summable
in trace norm.  Each individual term converges by
(MA-MO.1416), so dominated summation proves trace-norm
convergence of the series.  The same proof works when
$L_\epsilon\to L_0$ in Schwartz space.  This is the required
replacement for the stronger common-modulus estimate in MO030.

For $y\in\mathcal E$, the original wide inverse $F_{\omega,y}$
can first be shifted to every fixed $b>a$, because that line
has a fixed margin from all zeros.  Hence its actual identity is
$V_bF_{\omega,y}=A_\omega(b+i\,\cdot)^{-1}V_by$.
For fixed $y\in\mathcal E$, $V_by\to V_ay$ in Schwartz space
as $b\downarrow a$ on the indicated compact range of real parts.
The just-proved $H^{-r}$ convergence, also after multiplication
by every polynomial in frequency, identifies its boundary value:
\[
 V_aF_{\omega,y}=j_{\omega,+}V_ay,\qquad
 D_x^j(e^{ax}F_{\omega,y})\in
                   L^2((1+x^2)^{-r}\,dx)\quad(j\ge0).
 \tag{MA-MO.1419}
\]
The local $H^1$ estimate therefore gives polynomial left bounds
for this actual function and its derivatives.  Its original
right tail has every stronger exponential weight.  Thus
$qF_{\omega,y}$ is in the original local Schwartz-cup source,
and the boundary prescription agrees with the original maps.
There is no replacement of an original inverse by a chosen
unrelated boundary distribution.

The nonsingular factors in the actual removable commutator and
crossed kernels (MA-MO.718)--(MA-MO.719) converge in Schwartz
space as $b\downarrow a$.  This uses only forward-multiplier
bounds, the compact smooth test and cutoff derivative, and the
separated denominator $\Re(-z-w)=-2b\le-2a$.
Apply (MA-MO.1418) to their input reciprocals.  In the crossed
term the reciprocal is that of the dual sector, unchanged from
the original source formula.  The first row is exactly
\[
 \mathscr R_{\eta,a}^{\partial}(g)y
      =K_{\eta,a}^{\partial}(g)y-N_a(g)y
                         +O_{\eta,a}^{\partial}(g)Ry,
 \qquad N_a(g)=C_gM_\phi.
 \tag{MA-MO.1420}
\]
The half-line map $N_a(g):H_{\Sigma,a,s}\to H_{a,s}$ is
trace class by its physical kernel (MA-MO.332).
The other terms are trace class by (MA-MO.1418) and the
exact Fourier identification $H_{a,s}\simeq H^r$.
On each smooth input, the row is obtained by passing the
original $b>a$ formula to the limit, keeping that input and
its actual smooth theta output fixed.  No bounded embedding
between whole varying-width Hilbert spaces is assumed.
The reflected dual row agrees on the dense smooth common sector.
Its paired map therefore has range in the closed diagonal of
$H_{a,s}\oplus H_{-a,s}$, giving the trace-class endomorphism
in the exact inherited sum norm.

\emph{4. The rational polar completion and evaluated identities.}
The two rational cutoff moments are also controlled by
(MA-MO.1416).  With $p_a(x)=\phi(x)e^{(1/2-a)x}$,
\[
 \lambda_0(y)=\sum_n\left\langle d_{n,0}(u),
                 \widehat p_a(-u)(V_ay)(u)\right\rangle,
                \qquad\lambda_1(y)=\lambda_0(Ry).
 \tag{MA-MO.1421}
\]
The localized multiplier norms of the Schwartz function
$\widehat p_a(-u)$ decay rapidly; the sum of functional norms
is finite.  Passage from the original lines $b>a$, justified
above, identifies these with the two actual physical moments.
The left bound in (MA-MO.1419) and the wide right weights
make the full Mellin integral at $1/2$ converge.  The original
zeta inverse has a zero there, so $\mathcal MF_y(1/2)=0$.
Thus the actual source vector has the same two values
$h_y(0)=-\lambda_1(y)$ and $\int h_y=-\lambda_0(y)$.

Consequently the polar-subtracted diagonal rows and original
source columns assemble exactly as in (MA-MO.921)--(MA-MO.922):
\[
 \begin{aligned}
 \mathscr D_g(y,c)
     &=(T_g^\partial y,\nu_g(y))+\sum_{j=0}^1c_jC_gDs_j,\\
 \nu_g(y)&=(-G(-1/2)\lambda_1(y),-G(1/2)\lambda_0(y)).
 \end{aligned}
 \tag{MA-MO.1422}
\]
The first component includes the unchanged finite corrections
$G(-1/2)\lambda_1(y)e_0+G(1/2)\lambda_0(y)e_1$.
All terms are trace class or finite rank in the specified
source-section completion.  This proves the rational statement.

Every boundary multiplicity is at most the corresponding
$m_{\omega,n}<p$, and reflection supplies the other edge.
Hence $s>2m_\rho-1$ for every boundary zero.  Interior
multiplicities are automatically fully retained.  The actual
full operator, this full retention, and the ideal trace of
MA-MO-028 are exactly the hypotheses used in MA-MO-029.
Its relative-trace and determinant algebra therefore applies
without any finite-divisor assumption.  The source and primitive
Euler factors, both polar signs, and the $g(0)=0$ prime-formula
condition remain unchanged.
Finally (MA-MO.1418) involves only finitely many kernel seminorms,
which are controlled by finitely many derivatives of $g$ on a
fixed compact support.  This proves the asserted trace-norm
continuity.  Density of the original cores gives uniqueness.
\end{proof}

\paragraph{Concrete sufficient cases and exact limitations.}
MA-MO-030 is contained in this theorem: if $k_{\omega,n}\le M<p$,
then (MA-MO.1402) holds and every $B_{\omega,n}=1$, regardless
of inward distances or boundary collisions.  Beyond that class,
suppose there are no boundary zeros, $k_{\omega,n}\le M$,
and $d_{\omega,n}\ge c(2+|n|)^{-D_0}$.  Then (MA-MO.1405)
holds for every fixed $s\ge0$, without requiring $s>2M-1$.
Unbounded $k_{\omega,n}$ are also permitted whenever their
displayed depth budget is polynomial.  In the supercritical
case its logarithmic form is
\[
       (k_{\omega,n}-p)\log(1/d_{\omega,n})=O(\log(2+|n|)).
 \tag{MA-MO.1423}
\]
At the critical integer $k=p$ the separate square-root logarithm
in (MA-MO.1404) is essential.

The estimate uses a nonoscillatory bound for the whole boundary
subproduct.  If its total multiplicity is at least $p$, even at
distinct ordinates with each individual order below $p$, the
argument above does not apply.  This is a limitation of this
sufficient condition, not a proved operator obstruction in that
case.  Likewise failure of the positive-depth budget is not
asserted sufficient for failure of the actual operator: ordinate
separation and cancellations may improve the bound.
MA-MO-033 gives a distinct necessary estimate and a conditional
obstruction for a specified tight cluster configuration.
Neither result proves that any proposed configuration occurs.
The general no-margin full-retention arithmetic question remains
open; A or M supplies the previously separated fixed-gap branch.

\paragraph{Sources and trace-framework placement.}
The fixed-strip polynomial bound, Euler lower bound and ordinary
local zero count are the unconditional inputs in MT-M18.5--9.
The actual inverse, local source, Fourier phases, kernel lemma and
polar sections are the already proved MA-MO-018/022/027/028/030
interfaces.  The new estimates are (MA-MO.1409)--(MA-MO.1413)
and their blockwise dominated trace-norm assembly.  The source
boundary distribution is always tied back to the original wide
inverse before it is used in a completed operator.
In the inherited cote~1 framework, BR-TR-01/02, scans 079--080
and 111--112, finite-length $\mathbb C[z]$-module data give the
local multiplicities and exterior-power determinants.  The new
analytic budget and actual nuclear construction provide additional
structure; no positivity or arithmetic zero bound follows from
the finite-length $K_0$ classes alone.
